
\documentclass{article}
\usepackage{colortbl}
\usepackage{makecell}
\usepackage{multirow}
\usepackage{supertabular}

\begin{document}

\newcounter{utterance}

\centering \large Interaction Transcript for game `bbh', experiment `object\_counting', episode 154 with reflexion{-}llama{-}3.1{-}8b{-}instruct.
\vspace{24pt}

{ \footnotesize  \setcounter{utterance}{1}
\setlength{\tabcolsep}{0pt}
\begin{supertabular}{c@{$\;$}|p{.15\linewidth}@{}p{.15\linewidth}p{.15\linewidth}p{.15\linewidth}p{.15\linewidth}p{.15\linewidth}}
    \# & \multicolumn{2}{c}{Player} && \multicolumn{2}{c}{Game Master} \\
    \hline

    \theutterance \stepcounter{utterance}  
    & & & \multicolumn{4}{p{0.6\linewidth}}{
        \cellcolor[rgb]{0.9,0.9,0.9}{
            \makecell[{{p{\linewidth}}}]{
                \texttt{\tiny{[P1$\langle$GM]}}
                \texttt{You have attempted this same task before and did not succeed. These are the plans you wrote after those attempts.} \\
\texttt{Trial \#1: New plan: I failed to count the fruits in the last question because I didn't consider the possibility of a fruit being mentioned multiple times. In the question "I have an orange, a raspberry, two peaches, a blackberry, an apple, a grape, a nectarine, and three plums. How many fruits do I have?", I should have counted the peaches and plums as multiple fruits. Next time, I will make sure to count each fruit individually, even if it is mentioned multiple times.} \\
\\ 
\texttt{Trial \#2: New plan: I failed to count the fruits in the last question because I didn't consider the possibility of a fruit being mentioned multiple times. However, in the last question, the fruit was not mentioned multiple times, but rather, I failed to count the fruits correctly because I didn't consider the possibility of a fruit being mentioned at all. In the question "I have a banana, a nectarine, a raspberry, an apple, and a peach. How many fruits do I have?", I should have counted the fruits correctly without any issues. Next time, I will make sure to count each fruit individually, even if it is mentioned only once, and also consider the possibility of a fruit being mentioned multiple times in a future question.} \\
\\ 
\\ 
\texttt{Questions that involve enumerating objects and asking the model to count them.} \\
\\ 
\texttt{Q: I have a blackberry, a clarinet, a nectarine, a plum, a strawberry, a banana, a flute, an orange, and a violin. How many fruits do I have?} \\
\texttt{A: 6} \\
\\ 
\texttt{Q: I have an orange, a raspberry, two peaches, a blackberry, an apple, a grape, a nectarine, and three plums. How many fruits do I have?} \\
\texttt{A: 11} \\
\\ 
\texttt{Q: I have a lettuce head, a head of broccoli, an onion, a stalk of celery, two carrots, a garlic, and a yam. How many vegetables do I have?} \\
\texttt{A: 8} \\
\\ 
\texttt{Q: I have a banana, a nectarine, a raspberry, an apple, and a peach. How many fruits do I have?} \\
\texttt{A:} \\
            }
        }
    }
     \\ \\

    \theutterance \stepcounter{utterance}  
    & \multicolumn{4}{p{0.6\linewidth}}{
        \cellcolor[rgb]{0.561,0.737,0.561}{
            \makecell[{{p{\linewidth}}}]{
                \texttt{\tiny{[P1$\rangle$GM]}}
                \texttt{To answer the question "I have a banana, a nectarine, a raspberry, an apple, and a peach. How many fruits do I have?", I will follow the plan I wrote after my previous attempts.} \\
\\ 
\texttt{I will count each fruit individually, even if it is mentioned only once. In this case, I have the following fruits:} \\
\\ 
\texttt{1. Banana} \\
\texttt{2. Nectarine} \\
\texttt{3. Raspberry} \\
\texttt{4. Apple} \\
\texttt{5. Peach} \\
\\ 
\texttt{Since each fruit is counted individually, I will count them as separate fruits. Therefore, I have 5 fruits.} \\
            }
        }
    }
    & & \\ \\

    \theutterance \stepcounter{utterance}  
    & & & \multicolumn{2}{p{0.3\linewidth}}{
        \cellcolor[rgb]{0.9,0.9,0.9}{
            \makecell[{{p{\linewidth}}}]{
                \texttt{\tiny{[GM$|$GM]}}
                \texttt{Target: 5} \\
            }
        }
    }
    & & \\ \\

    \theutterance \stepcounter{utterance}  
    & & & \multicolumn{2}{p{0.3\linewidth}}{
        \cellcolor[rgb]{0.9,0.9,0.9}{
            \makecell[{{p{\linewidth}}}]{
                \texttt{\tiny{[GM$|$GM]}}
                \texttt{game\_result = LOSE} \\
            }
        }
    }
    & & \\ \\

    \theutterance \stepcounter{utterance}  
    & & \multicolumn{4}{p{0.6\linewidth}}{
        \cellcolor[rgb]{0.95,0.95,0.95}{
            \makecell[{{p{\linewidth}}}]{
                \texttt{\tiny{[P1$\leftrightarrow$P]}}
                \texttt{Post{-}episode Reflexion (not sent to GM):} \\
\texttt{New plan: I failed to count the fruits in the last question because I didn't consider the possibility of a fruit being mentioned only once, and also because I didn't consider the possibility of a fruit being mentioned multiple times. However, in the last question, the fruit was not mentioned multiple times, but rather, I failed to count the fruits correctly because I didn't consider the possibility of a fruit being mentioned at all. In the question "I have a banana, a nectarine, a raspberry, an apple, and a peach. How many fruits do I have?", I should have counted the fruits correctly without any issues. Next time, I will make sure to count each fruit individually, even if it is mentioned only once, and also consider the possibility of a fruit being mentioned multiple times in a future question. To achieve this, I will create a list of all the fruits mentioned in the question and then count the number of items in the list. This way, I will ensure that I count each fruit individually, regardless of whether it is mentioned multiple times or not.} \\
            }
        }
    }
    & \\ \\

\end{supertabular}
}

\end{document}
